% Préambule commun des schémas TikZ — charte Liouaï
% Chaque schéma commence par % Préambule commun des schémas TikZ — charte Liouaï
% Chaque schéma commence par % Préambule commun des schémas TikZ — charte Liouaï
% Chaque schéma commence par \input{_preambule} puis \begin{document}.
\documentclass[tikz,border=6pt]{standalone}
\usepackage[T1]{fontenc}
\usepackage[utf8]{inputenc}
\usepackage[scaled=0.95]{helvet}
\renewcommand{\familydefault}{\sfdefault}
\usepackage{amsmath,amssymb}
\usepackage[european]{circuitikz}
\usetikzlibrary{arrows.meta,positioning,calc,angles,quotes,decorations.pathreplacing,backgrounds}

% Couleurs du logo Liouaï Digital Education
\definecolor{navy}{HTML}{243743}
\definecolor{vert}{HTML}{28B78D}
\definecolor{vertfonce}{HTML}{167A5C}
\definecolor{vertclair}{HTML}{EAF7F2}
\definecolor{gris}{HTML}{A1A9AE}
\definecolor{gristexte}{HTML}{5C6870}
\definecolor{fond}{HTML}{F3F5F6}
\definecolor{jaune}{HTML}{E5C800}
\definecolor{jauneclair}{HTML}{FFFBD6}

\tikzset{
  every picture/.style={line cap=round, line join=round, color=navy, text=navy},
  fleche/.style={-{Stealth[length=2.6mm,width=2mm]}, line width=1.1pt},
  vecteur/.style={fleche, line width=1.4pt},
  titre/.style={font=\bfseries\small, text=navy},
  note/.style={font=\footnotesize, text=gristexte},
  boite/.style={draw=navy, line width=0.9pt, rounded corners=5pt, fill=fond, align=center, inner sep=6pt},
}
 puis \begin{document}.
\documentclass[tikz,border=6pt]{standalone}
\usepackage[T1]{fontenc}
\usepackage[utf8]{inputenc}
\usepackage[scaled=0.95]{helvet}
\renewcommand{\familydefault}{\sfdefault}
\usepackage{amsmath,amssymb}
\usepackage[european]{circuitikz}
\usetikzlibrary{arrows.meta,positioning,calc,angles,quotes,decorations.pathreplacing,backgrounds}

% Couleurs du logo Liouaï Digital Education
\definecolor{navy}{HTML}{243743}
\definecolor{vert}{HTML}{28B78D}
\definecolor{vertfonce}{HTML}{167A5C}
\definecolor{vertclair}{HTML}{EAF7F2}
\definecolor{gris}{HTML}{A1A9AE}
\definecolor{gristexte}{HTML}{5C6870}
\definecolor{fond}{HTML}{F3F5F6}
\definecolor{jaune}{HTML}{E5C800}
\definecolor{jauneclair}{HTML}{FFFBD6}

\tikzset{
  every picture/.style={line cap=round, line join=round, color=navy, text=navy},
  fleche/.style={-{Stealth[length=2.6mm,width=2mm]}, line width=1.1pt},
  vecteur/.style={fleche, line width=1.4pt},
  titre/.style={font=\bfseries\small, text=navy},
  note/.style={font=\footnotesize, text=gristexte},
  boite/.style={draw=navy, line width=0.9pt, rounded corners=5pt, fill=fond, align=center, inner sep=6pt},
}
 puis \begin{document}.
\documentclass[tikz,border=6pt]{standalone}
\usepackage[T1]{fontenc}
\usepackage[utf8]{inputenc}
\usepackage[scaled=0.95]{helvet}
\renewcommand{\familydefault}{\sfdefault}
\usepackage{amsmath,amssymb}
\usepackage[european]{circuitikz}
\usetikzlibrary{arrows.meta,positioning,calc,angles,quotes,decorations.pathreplacing,backgrounds}

% Couleurs du logo Liouaï Digital Education
\definecolor{navy}{HTML}{243743}
\definecolor{vert}{HTML}{28B78D}
\definecolor{vertfonce}{HTML}{167A5C}
\definecolor{vertclair}{HTML}{EAF7F2}
\definecolor{gris}{HTML}{A1A9AE}
\definecolor{gristexte}{HTML}{5C6870}
\definecolor{fond}{HTML}{F3F5F6}
\definecolor{jaune}{HTML}{E5C800}
\definecolor{jauneclair}{HTML}{FFFBD6}

\tikzset{
  every picture/.style={line cap=round, line join=round, color=navy, text=navy},
  fleche/.style={-{Stealth[length=2.6mm,width=2mm]}, line width=1.1pt},
  vecteur/.style={fleche, line width=1.4pt},
  titre/.style={font=\bfseries\small, text=navy},
  note/.style={font=\footnotesize, text=gristexte},
  boite/.style={draw=navy, line width=0.9pt, rounded corners=5pt, fill=fond, align=center, inner sep=6pt},
}

\begin{document}
\begin{tikzpicture}[scale=0.95]
\def\panneau#1#2#3{\node[titre] at (#1,2.15) {#2}; \node[note] at (#1,-2.0) {#3};}
% 1. départ
\begin{scope}
\draw[gris, line width=0.5pt] (-1.5,0) -- (1.5,0); \draw[gris, line width=0.5pt] (0,-1.5) -- (0,1.5);
\draw[navy, line width=1pt] (0,0) circle (1);
\draw[vecteur, vertfonce] (0,0) -- (40:1) node[above right=-2pt] {$\mathbf{v}_1$};
\draw[vecteur, jaune!85!navy] (0,0) -- (130:1) node[above left=-2pt] {$\mathbf{v}_2$};
\end{scope}
\panneau{0}{départ}{cercle unité}
\draw[fleche, navy] (1.7,0) -- (2.7,0) node[midway, above, font=\small] {$V^{\mathsf{T}}$};
% 2. rotation
\begin{scope}[xshift=4.4cm]
\draw[gris, line width=0.5pt] (-1.5,0) -- (1.5,0); \draw[gris, line width=0.5pt] (0,-1.5) -- (0,1.5);
\draw[navy, line width=1pt] (0,0) circle (1);
\draw[vecteur, vertfonce] (0,0) -- (1,0) node[below right=-2pt] {$\mathbf{e}_1$};
\draw[vecteur, jaune!85!navy] (0,0) -- (0,1) node[above left=-2pt] {$\mathbf{e}_2$};
\end{scope}
\panneau{4.4}{1. tourner}{$\mathbf{v}_1, \mathbf{v}_2$ viennent sur les axes}
\draw[fleche, navy] (6.1,0) -- (7.1,0) node[midway, above, font=\small] {$\Sigma$};
% 3. étirement
\begin{scope}[xshift=9.2cm]
\draw[gris, line width=0.5pt] (-1.9,0) -- (1.9,0); \draw[gris, line width=0.5pt] (0,-1.5) -- (0,1.5);
\draw[navy, line width=1pt] (0,0) ellipse (1.7 and 0.55);
\draw[vecteur, vertfonce] (0,0) -- (1.7,0) node[above=1pt] {$\sigma_1$};
\draw[vecteur, jaune!85!navy] (0,0) -- (0,0.55) node[above left=-2pt] {$\sigma_2$};
\end{scope}
\panneau{9.2}{2. étirer}{selon les axes, par $\sigma_1 \geq \sigma_2 \geq 0$}
\draw[fleche, navy] (11.2,0) -- (12.2,0) node[midway, above, font=\small] {$U$};
% 4. rotation finale
\begin{scope}[xshift=14.3cm]
\draw[gris, line width=0.5pt] (-1.9,0) -- (1.9,0); \draw[gris, line width=0.5pt] (0,-1.5) -- (0,1.5);
\draw[navy, line width=1pt, rotate=25] (0,0) ellipse (1.7 and 0.55);
\draw[vecteur, vertfonce] (0,0) -- (25:1.7) node[above=1pt] {$\sigma_1\mathbf{u}_1$};
\draw[vecteur, jaune!85!navy] (0,0) -- (115:0.55) node[left=1pt] {$\sigma_2\mathbf{u}_2$};
\end{scope}
\panneau{14.3}{3. tourner}{résultat : l'ellipse de $A$}
\node[boite, font=\small] at (7.2,-3.3) {$A = U\,\Sigma\,V^{\mathsf{T}}$ \qquad toute matrice = rotation, étirement selon des axes perpendiculaires, rotation};
\end{tikzpicture}
\end{document}
